%%%%%%%%%%%%
%% Please rename this main.tex file and the output PDF to
%% [lastname_firstname_graduationyear]
%% before submission.
%%
%% This .tex file is for use with BibLaTeX. Please use
%% main-bibtex.tex instead if you prefer BibTeX.
%%%%%%%%%%%%

\documentclass[12pt]{caltech_thesis}
\usepackage[hyphens]{url}
\usepackage{lipsum}
\usepackage{graphicx}

\usepackage{todonotes}

%% Tentative: newtx for better-looking Times
\usepackage[utf8]{inputenc}
\usepackage[T1]{fontenc}
\usepackage{newtxtext,newtxmath}

% Must use biblatex to produce the Published Contents and Contributions, per-chapter bibliography (if desired), etc.
\usepackage[
    backend=biber,natbib,
    % IMPORTANT: load a style suitable for your discipline
    style=authoryear
]{biblatex}

% Name of your .bib file(s)
\addbibresource{example.bib}
\addbibresource{ownpubs.bib}

\begin{document}

% Do remember to remove the square bracket!
\title{[Thesis Title]}
\author{[Your Full Name]}

\degreeaward{[Name of Degree]}                 % Degree to be awarded
\university{California Institute of Technology}    % Institution name
\address{Pasadena, California}                     % Institution address
\unilogo{caltech.png}                                 % Institution logo
\copyyear{[Year Degree Conferred]}  % Year (of graduation) on diploma
\defenddate{[Exact Date]}          % Date of defense

\orcid{[Author ORCID]}

%% IMPORTANT: Select ONE of the rights statement below.
\rightsstatement{All rights reserved\todo[size=\footnotesize]{Choose one from the choices in the source code!! And delete this \texttt{todo} when you're done that. :-)}}
% \rightsstatement{All rights reserved except where otherwise noted}
% \rightsstatement{Some rights reserved. This thesis is distributed under a [name license, e.g., ``Creative Commons Attribution-NonCommercial-ShareAlike License'']}

%%  If you'd like to remove the Caltech logo from your title page, simply remove the "[logo]" text from the maketitle command
\maketitle[logo]
%\maketitle

\begin{acknowledgements} 	 
   [Add acknowledgements here. If you do not wish to add any to your thesis, you may simply add a blank titled Acknowledgements page.]
\end{acknowledgements}

\begin{abstract}
   [This abstract must provide a succinct and informative condensation of your work. Candidates are welcome to prepare a lengthier abstract for inclusion in the dissertation, and provide a shorter one in the CaltechTHESIS record.]
\end{abstract}

%% Uncomment the `iknowhattodo' option to dismiss the instruction in the PDF.
\begin{publishedcontent}%[iknowwhattodo]
% List your publications and contributions here.
\nocite{Cahn:etal:2015,Cahn:etal:2016}
\end{publishedcontent}

\tableofcontents
\listoffigures
\listoftables
\printnomenclature

\mainmatter

\chapter{Introduction}
Start off all chapters with \verb|chapter|. \index{chapter!numbered} \verb|\extrachapter| will give you an unnumbered chapter that's added to the Table of Contents. \index{chapter!unnumbered}

Here's an example of a citation \citep{GMP81}. Here's another \citep{PP98}. These will appear in the big bibliography at the end of the thesis.
\index{bibliography}

If you're new to \LaTeX{} and would like to begin by learning the basics, please see our free online course available at:\\ \url{https://www.overleaf.com/latex/learn/free-online-introduction-to-latex-part-1} \index{LaTeX@\LaTeX}

You can define nomenclatures \index{nomenclature} as you talk about key terms in your thesis. So what's a galaxy? \nomenclature{Galaxy}{A system of stars independent from all other systems}


\section{$G_0W_0$}
\subsection{Interpretive procedure}
The procedure that was used to compute the quasiparticle energies is given by the below equation:
\begin{equation}
    \delta_{pq}F_{pq}^{HF}[\gamma^{DFT }] + \Sigma_{p}^{corr}(\varepsilon_{p}^{QP}) = \varepsilon_{p}^{QP}
\end{equation}
We explain the notation starting from left to right. The first term response to taking the diagonal $\delta _{pq}$ of the harte fog fog matrix $F_{pq}^{HF}$ evaluated at a given electron density $\gamma $. These electron densities are obtained form a previous mean field calculation, either $\gamma _{DFT}$ or $\gamma _{HF}$. 
\begin{figure}[hbt!]
\centering
\includegraphics[width=.3\textwidth]{caltech.png}
\caption{This is a figure}\label{fig:logo}
\index{figures}
\end{figure}

\subsection{This is a subsection}

\begin{table}[hbt!]
\centering
\begin{tabular}{ll}
\hline
Area & Count\\
\hline
North & 100\\
South & 200\\
East & 80\\
West & 140\\
\hline
\end{tabular}
\caption{This is a table}\label{tab:sample}
\index{tables}
\end{table}

\lipsum[3] \nomenclature{Asteroid}{A very small planet ranging from 1,000 km to less than one km in diameter. Asteroids are found commonly around other larger planets}

\lipsum[4-5] 

Here's an endnote.\endnote{Endnotes are notes that you can use to explain text in a document.}

\section{This is Another Section}
\lipsum[6-7] 

\chapter{$G_0W_0$}
\section{Interpretive procedure}
The procedure that was used to compute the quasiparticle energies is given by the below equation:
\begin{equation}
    \delta_{pq}F_{pq}^{HF}[\gamma^{DFT}] + \Sigma_{p}^{corr}(\varepsilon_{p}^{QP}) = \varepsilon_{p}^{QP}
\end{equation}
We explain the notation starting from left to right. The first term responds to taking the diagonal $\delta_{pq}$ of the Hartree-Fock matrix $F_{pq}^{HF}$ evaluated at a given electron density $\gamma$. These electron densities are obtained from a previous mean-field calculation, either $\gamma_{DFT}$ or $\gamma_{HF}$. 
\subsection{The Fock Matrix}
In the basis of atomic orbitals, this is given by:
\begin{equation}
F_{\mu\nu}^{HF} = h_{\mu\nu} + \sum_{\lambda\sigma}P_{\lambda\sigma}(\mu\nu|\lambda\sigma) - \frac{1}{2}\sum_{\lambda\sigma}P_{\lambda\sigma}(\mu\lambda|\nu\sigma)
\end{equation}
where $h_{\mu\nu}$ is the one-electron part of the Hamiltonian, $P_{\lambda\sigma}$ is the density matrix, and $(\mu\nu|\lambda\sigma)$ is one of the two-electron integrals. Cite Szabo. This is the simple form of the Hartree-Fock matrix that we want to use here and not the DFT Fock matrix. We transform this Fock matrix into the MO basis with:
\begin{equation}
   F_{pq} = \sum_{\mu} \sum_{\nu} C_{\mu p}^{*}F_{\mu\nu}C_{\nu q}
\end{equation}
where $C$ is the matrix of MO coefficients. Another useful identity is for the density matrix in terms of the MO coefficients from the mean-field calculation: 
\begin{equation}
P_{\mu\nu} = 2\sum_{i=1}^{N/2}C_{\mu i}C_{\nu i}^{*}
\end{equation}
We note that the sum runs only over the $N/2$ occupied \emph{spatial} orbitals.
\subsection{Real Correlation-Solve Energy}
This is the second term in 2.1. It is dynamic, as opposed to the previous Fock term that was discussed, as it is updated with a new quasiparticle energy in each iteration. In the case of the $G_0W_0$ approximation, we are only interested in the diagonal element of $\Sigma^{corr}$ corresponding to the orbital with index $p$. This function is evaluated at the QP energy $\varepsilon_{p}^{QP}$ just obtained in the previous iteration. We will go into greater detail about the form of $\Sigma^{corr}$ in the next chapter.


\chapter{This is the Third Chapter}

\publishedas{Cahn:etal:2015}

[You can have chapters that w3ere published as part of your thesis. The text style of the body should be single column, as it was submitted to the publisher, not formatted as the publisher did.]

\chapter{This is the Fourth Chapter}
\chapter{This is the Fifth Chapter}
\chapter{This is the Sixth Chapter}
\chapter{This is the Seventh Chapter}
\chapter{This is the Eighth Chapter}

\printbibliography[heading=bibintoc]

\appendix

\chapter{Questionnaire}
\chapter{Consent Form}

\printindex

\theendnotes

%% Pocket materials at the VERY END of thesis
\pocketmaterial
\extrachapter{Pocket Material: Map of Case Study Solar Systems} 


\end{document}
